\subsection{THETA3-C18: tame theta ratios have an integral common frame}\label{C18}

Let $p\geq5$, let $R$ be the ring of integers of a finite extension of
$\Q_p$, and let $\mathcal E/R$ be the smooth elliptic curve
$y^2=4x^3-g_2x-g_3$ with $g_2,g_3\in R$ and unit discriminant.
Use its integral parameter $t=-2x/y$ and its formal logarithm
$\lambda(t)$ normalized by the invariant differential $dx/y$.
The letter $\mathcal E$ here denotes a curve, not a CM field.
Enlarge $R$ by a finite extension whenever the torsion points below
must be made rational.

For a compatible system of division points $\widetilde T=(T_n)$,
suppose $T_1$ is a nonzero torsion point of order $N$ prime to $p$.
Write $s_T(t)$ for its algebraic sigma series in
Katagiri--Kobayashi, \emph{p-adic properties of division polynomials
and algebraic sigma functions}, arXiv:2609.07024v1,
Definition~3.1, p.~9. Their Remark~3.2 on that page identifies it over
$\C$ with
\[
 s_T(t)=\sigma_{z_T}(\lambda(t)),\qquad
 \sigma_{z_T}(z)=
 e^{-\eta(z_T)(z+z_T/2)}\sigma(z+z_T).
\]
The torsion lift, differential, and normalization are part of this
notation. Theorem~2.7, pp.~7--8, supplies the algebraic value and
Taylor-coefficient statement. For an arbitrary compatible system,
one can choose a rational lattice lift realizing the finite datum
$T_{4N}$ on which these constructions depend.
The following is a precise integral consequence of their
construction, rather than an integrality assertion for the raw sigma
series.

\paragraph{Common-factor statement.}
There are units $u_T(t)\in R[[t]]^\times$ such that
\begin{equation}\label{C18-factor}
 s_T(t)=q(t)u_T(t),\qquad
 q(t)=\frac{\sigma(\lambda(t))}{t},\qquad
 u_T(0)=\sigma_{\rm alg}(\widetilde T)\in R^\times.
\end{equation}
The same $q(t)\in\operatorname{Frac}(R)[[t]]^\times$ works for every
such tame torsion point. Consequently, for two lists of the same
length,
\begin{equation}\label{C18-ratio}
 R_{\mathbf T,\mathbf U}(t):=
 \frac{\prod_{i=1}^r s_{T_i}(t)}
      {\prod_{i=1}^r s_{U_i}(t)}
       =\frac{\prod_i u_{T_i}(t)}{\prod_i u_{U_i}(t)}
       \in R[[t]]^\times.
\end{equation}
After pullback by a common prime-to-$p$ multiplication isogeny,
this is the completion of a rational function on
$\mathcal E$ that is a unit along the identity residue disc.
On that cover it is the ratio of sections of equal tensor powers
of the same line bundle; before pullback a torsion flat twist can
remain. This integral-unit conclusion is specific to equal tensor
degrees and prime-to-$p$ torsion. A weaker derivative estimate for
arbitrary integer tensor degree is proved below.

\paragraph{Proof, including the integral model.}
The rational function $g_T$ used in Definition~3.1 has
\[
 \operatorname{div}(g_T)=[N]^*(\tau_{T_1}^*(O)-(O)),\qquad
 (t g_T)(O)=N^{-1}\sigma_{\rm alg}(\widetilde T).
\]
By their Lemma~3.3(1), pp.~9--10, the latter sigma value is a
unit. Here is the resulting local integral calculation.
Because $N$ is prime to $p$, all the torsion sections in this
divisor are disjoint over $R$ after the finite extension. The
only one through the special-fiber origin is the simple pole
along $O$. Any vertical contribution to the divisor of $g_T$
is an integer multiple of the smooth irreducible special fiber.
In the regular local ring at its origin, cancellation of the
horizontal pole gives $t g_T=\varpi^a v$ with $v$ a unit and
$a\in\Z$. Its restriction to $O$ is a unit by the displayed
normalization. Hence $a=0$, and
\[
 h_T(t):=t g_T(t)\in R[[t]]^\times.
\]
This argument can equivalently be made directly with the
horizontal rational function in the proof of Lemma~3.3.

Multiplication by $N$ on the integral formal group has expansion
$[N](t)=Nt+\cdots$ with $N\in R^\times$. Its compositional inverse
$a_N(t)=[N^{-1}](t)$ belongs to $tR[[t]]$, with linear coefficient
$N^{-1}$. Therefore
\[
 u_T(t):=t g_T(a_N(t))
       =\frac{t}{a_N(t)}h_T(a_N(t))\in R[[t]]^\times,
 \qquad u_T(0)=N\,h_T(0)=\sigma_{\rm alg}(\widetilde T).
\]
Definition~3.1 now gives \eqref{C18-factor}, and cancellation
gives \eqref{C18-ratio}.

To verify the rational-function assertion, take a common multiple
$M$ of all the orders $N_i,N_i'$ in the two lists, still prime
to $p$. The formal group identity
$a_{N_i}([M](t))=[M/N_i](t)$ gives directly
\[
 R_{\mathbf T,\mathbf U}([M](t))=
 \frac{\prod_i[M/N_i]^*g_{T_i}(t)}
      {\prod_i[M/N_i']^*g_{U_i}(t)}.
\]
The common sigma factors cancel before this substitution.
The right side is an explicit rational function on $\mathcal E$
over the same finite torsion extension. Its completion is an
integral unit by \eqref{C18-ratio} and the integral multiplication
series. This proves the asserted rationality on the finite \'etale
isogeny cover without asserting descent to the original curve.
\hfill$\square$

\paragraph{Derivative consequence.}
Suppose the integral formal parameter has been identified by a
unit-linear integral change of coordinates with the exact height-two
model of C11, $[p](T)=pT+\alpha T^{p^2}$ with $\alpha\in R^\times$.
For its normalized invariant derivation $D$, C11 applies to
\eqref{C18-ratio} and gives, for every $a\geq1$,
\[
 v_p\bigl(D^a R_{\mathbf T,\mathbf U}(0)\bigr)
       \geq\left\lfloor\frac{ap}{p^2-1}\right\rfloor-1.
\]
The assertion survives multiplication by an integral unit and
unit-linear integral changes of parameter, with the invariant
differential normalization retained. This is an actual theta
quotient to which the integral-function estimate applies.

For a single translated theta section the same calculation gives
the exact connection conversion
\[
 q^{-1}D(qf)=Df+\frac{Dq}{q}f.
\]
It separates the integral function $u_T$ from the common connection
term $Dq/q$. The coefficient example in C11 shows why $q$ cannot
simply be declared integral. For a tensor quotient of total degree
zero this common factor cancels as in \eqref{C18-ratio}; for a
tensor of degree $d$ it remains as $q^d$. Thus the remaining
comparison with Alp\"oge's harmonic covariant theta derivatives
is the specified connection and its degree, rather than an
unspecified raw-theta integrality hypothesis. This statement does
not identify those archive derivatives with the degree-zero
quotients without their actual connection data.

\paragraph{The derivative of the common connection is integral and rational.}
Put $A=Dq/q$, with $D$ normalized by $dx/y$.  The classical identity
$(\log\sigma)''=-\wp$ and $x=\wp(z)$ give the exact identity
\begin{equation}\label{C18-connection-rational}
 B:=DA=D^2\log q=-x-D^2\log t
   =\frac{(Dt)^2-tD^2t-t^2x}{t^2}.
\end{equation}
This is a rational function on the elliptic curve, and
$B\in t^2R[[t]]$.  Here is a direct integral verification, which does
not assume that $q$ or $A$ is integral.  Put $u=t^2x$.  The Weierstrass
equation and $t=-2x/y$ give
\[
 4u^2(u-1)=g_2u t^4+g_3t^6.
\]
The implicit-function derivative at $(u,t)=(1,0)$ is $4\in R^\times$.
Thus $u\in1+t^4R[[t]]$, so $x=t^{-2}+t^2R[[t]]$.
Since $Dx=y$ and $Dy=6x^2-g_2/2$, direct differentiation gives
\[
 Dt=\frac{4x^3+g_2x+2g_3}{4x^3-g_2x-g_3}
    =1+\frac{g_2t^4}{2u}+\frac{3g_3t^6}{4u^2}
       \in1+t^4R[[t]].
\]
Consequently $D=(Dt)\,d/dt$ preserves $R[[t]]$ and
$D^2t\in t^3R[[t]]$.  The numerator in
\eqref{C18-connection-rational} belongs to $t^4R[[t]]$, proving the
claim.  An explicit expression, useful independently of the parameter,
is
\begin{align*}
 B&=9x-\frac{g_2}{2x}-\frac{g_3}{x^2}
       -\frac{(6x^2-g_2/2)^2}{y^2}\\
  &=\frac{-5g_2x^4-13g_3x^3+(g_2^2/4)x^2
                +(3g_2g_3/2)x+g_3^2}{x^2y^2}.
\end{align*}
Since $q\in1+t^4\operatorname{Frac}(R)[[t]]$, both
$\log q(0)$ and $D\log q(0)$ are zero.  Under the exact height-two
coordinate hypothesis already stated, C11 now gives
\[
 v_p\bigl(D^a\log q(0)\bigr)
   \geq\left\lfloor\frac{(a-2)p}{p^2-1}\right\rfloor-1
       \qquad(a\geq3),
\]
while $D^2\log q(0)=B(0)=0$.
The integral rational derivative need not have an integral primitive.
For example, C11's curve with $g_2=4,g_3=0$ gives
\[
 A=-\frac53t^3+\frac{11}{7}t^7+O(t^{11}),\qquad
 B=-5t^2+t^6+O(t^{10}).
\]
Thus at $p=7$ the connection term itself is not integral in this frame.

\paragraph{A weaker positive slope for every tame tensor degree.}
Retain the exact height-two coordinate hypothesis, so that the
operator congruence from C11 holds:
\[
 D^p\bigl(R[[t]]\bigr)\subset pR[[t]].
\]
For every integer $d$, every $v(t)\in R[[t]]$, and every $n\geq1$,
the series $f=q^dv$ satisfies
\begin{equation}\label{C18-weaker-slope}
 v_p\bigl(D^nf(0)\bigr)
 \geq\max\left\{0,\ v_p(n!)-\frac{np}{p^2-1}\right\}
 =\max\left\{0,\frac{n}{p^2-1}
                         -\frac{s_p(n)}{p-1}\right\}.
\end{equation}
Here $s_p(n)$ is the sum of the base-$p$ digits of $n$.
This applies in particular to every finite product or quotient of the
tame translated sigma series: its integer tensor degree is $d$, and
the remaining factor $v$ is an integral unit.  Degree zero still has
the stronger C11 estimate with slope $p/(p^2-1)$.
The slope in \eqref{C18-weaker-slope} is $1/(p^2-1)$; it does not assert
the stronger slope for a single theta section.

\emph{Proof.}
The operator congruence implies
$D^k(R[[t]])\subset p^{\lfloor k/p\rfloor}R[[t]]$ by iteration.
Let $l_j=D^j\log q(0)$.  In the additive coordinate $z=\lambda(t)$,
both $t(z)$ and $\sigma(z)$ are odd.  Hence $q(t(z))$ and its formal
logarithm are even.  As the logarithm starts in degree four,
$l_j=0$ for odd $j$ and for $j<4$.  For even $j\geq4$,
\[
 l_j=D^{j-2}B(0),\qquad
 v_p(l_j)\geq\left\lfloor\frac{j-2}{p}\right\rfloor.
\]
For every such $j$ one has the elementary inequality
\begin{equation}\label{C18-Bell-weight}
 v_p(j!)-\left\lfloor\frac{j-2}{p}\right\rfloor
          +\frac1{p-1}\leq\frac{jp}{p^2-1}.
\end{equation}
For $j\leq p-1$ this is immediate; for $j=p+1$ it is an equality.
The remaining even integers satisfy $j\geq p+3$.  Using
$v_p(j!)\leq(j-1)/(p-1)$ and
$\lfloor(j-2)/p\rfloor\geq(j-p-1)/p$, the left side is at most
$j/(p(p-1))+(p+1)/p$.  Its difference from the right side is
nonpositive because already at $j=p+3$ the right side minus this
upper bound is
$(p^2-3p-2)/(p(p^2-1))>0$; the difference then increases with $j$.

Bell's formula for the exponential gives
\[
 D^n(q^d)(0)=
 n!\!\sum_{\substack{m_j\geq0,\ \sum jm_j=n\\j\geq4\ {\rm even}}}
       \prod_j\frac{(d l_j)^{m_j}}{m_j!(j!)^{m_j}}.
\]
It is also a polynomial with integer coefficients in the integral
quantities $d l_j$, so these derivatives are integral.  In each term,
use $v_p(m_j!)\leq m_j/(p-1)$ and
\eqref{C18-Bell-weight}.  The resulting valuation is at least
$v_p(n!)-np/(p^2-1)$.  This proves the stated bound for $q^d$,
including negative $d$; $d=0$ is immediate.

Finally write the additive-coordinate series for $q^d$ and $v$ in
exponential Taylor form.  The coefficient of $z^k$ in $q^d$ has
valuation at least $-kp/(p^2-1)$.  The coefficient of $z^a$ in $v$
has valuation at least
\[
 \left\lfloor\frac ap\right\rfloor-v_p(a!)
  =-\sum_{r\geq2}\left\lfloor\frac a{p^r}\right\rfloor
  \geq-\frac{a}{p(p-1)}
  \geq-\frac{ap}{p^2-1}.
\]
Convolution gives the same coefficient bound for $q^dv$, and
multiplication by $n!$ gives the factorial lower bound in
\eqref{C18-weaker-slope}.  Integrality of all its derivatives follows
separately from Leibniz's formula, because the derivatives of $q^d$
and $v$ are integral.  Legendre's factorial formula gives the final
digit-sum expression. \hfill$\square$

\framework{This is an integral trivialization and specialization
calculation for a line-bundle ratio on a smooth proper model.
Its geometric setting is EGA~IV$_2$, Theorem~6.9.1, p.~153
(flat spreading after localization), and SGA~4$\frac12$,
\emph{Cycle}, Definition~2.3.2, p.~145, and Theorem~2.3.8,
p.~148 (cycle classes and pullback). The connection normalization
is additional arithmetic data, not a change to Grothendieck's
cycle-class or polarization statements.}
